\documentclass[11pt]{article}
\usepackage[a4paper,margin=1in]{geometry}
\usepackage{amsmath,amssymb,mathtools,bm}
\usepackage{enumitem,booktabs,array}
\usepackage{tcolorbox}
\usepackage{hyperref}
\usepackage[protrusion=true,expansion=false]{microtype}
\hypersetup{colorlinks=true,linkcolor=blue,urlcolor=blue,citecolor=blue}
\setlist[itemize]{topsep=2pt,itemsep=2pt}
\setlist[enumerate]{topsep=2pt,itemsep=2pt}
\newtcolorbox{keybox}{colback=blue!3!white,colframe=blue!40!black,boxrule=0.4pt,arc=1mm}
\newtcolorbox{alertbox}{colback=red!3!white,colframe=red!40!black,boxrule=0.4pt,arc=1mm}
\newtcolorbox{answerbox}{colback=green!3!white,colframe=green!40!black,boxrule=0.4pt,arc=1mm}
\newtcolorbox{drillbox}{colback=yellow!5!white,colframe=orange!60!black,boxrule=0.4pt,arc=1mm}
\newcommand{\EV}{\mathbb{E}}
\newcommand{\Var}{\mathrm{Var}}
\newcommand{\Cov}{\mathrm{Cov}}
\newcommand{\blank}[1]{\underline{\hspace{#1}}}

\title{W11-04: Fernando, Sharfman \& Uysal (2017, JFQA) \\
\large ``Corporate Environmental Policy and Shareholder Value: Following the Smart Money'' --- Active Recall Workbook}
\author{Banking \& Risk Management --- Exam Prep}
\date{\today}
\begin{document}\maketitle

\begin{keybox}
\textbf{Paper in one sentence.} FSU show that institutional investors overweight firms with \emph{intermediate} environmental performance: they avoid environmental ``laggards'' (due to risk) and also under-weight ``leaders'' (whose environmental investment may be inefficient), consistent with institutions pricing environmental risk without endorsing over-investment.

\textbf{Placement.} Distinguishes \emph{risk management} from \emph{moral preference} in institutional ESG allocation; a key counter to simple ``institutions love green'' narratives. Pairs with Chava (W11-03) and Sharfman-Fernando (W11-01).
\end{keybox}

\section*{Round 1: Complete Walk-Through}

\subsection*{1.1 Data and measures}
KLD environmental strengths and concerns for 1992--2008. Institutional holdings from Thomson 13-F. Define three bins: laggards (concerns $>$ strengths), neutral, leaders (strengths $>$ concerns).

\subsection*{1.2 Central test}
Compare institutional ownership across the three bins, controlling for firm characteristics.
\[
\text{InstOwn}_{it}=\alpha+\beta_L\cdot\text{Laggard}_{it}+\beta_H\cdot\text{Leader}_{it}+X_{it}'\gamma+\varepsilon_{it}.
\]
Reference is the neutral bin.

\subsection*{1.3 Headline results}
\begin{itemize}
\item $\beta_L < 0$: laggards are significantly \emph{underweighted} by institutions.
\item $\beta_H < 0$ (milder): leaders are also underweighted relative to neutral.
\item Effect stronger for pensions and long-horizon institutions than hedge funds.
\item ``Smart money'' reaction: institutions treat environmental risk as a cost but do not pay a premium for best-in-class green performance.
\end{itemize}

\subsection*{1.4 Interpretation}
\begin{itemize}
\item Risk-management motive: avoid environmental tail risk (laggards) $\Rightarrow$ lower holdings.
\item Efficiency motive: over-investment in environmental projects may destroy value; institutions reduce exposure.
\item Consistent with ``risk premium, not ethical premium'' reading of institutional ESG behaviour.
\end{itemize}

\subsection*{1.5 Caveats}
\begin{alertbox}
\begin{itemize}
\item KLD binning is coarse; measurement of ``leader'' vs.\ ``neutral'' may be noisy.
\item Results may have shifted post-2010 as ESG investing grew.
\item Endogeneity: firms target their environmental profile to institutional preferences.
\end{itemize}
\end{alertbox}

\section*{Round 2: Level 1 Fill-in}
\begin{enumerate}
\item Data: KLD strengths minus concerns, \blank{1cm}--\blank{1cm}.
\item Holdings source: Thomson \blank{1cm}-F.
\item Bins: laggards / neutral / \blank{1cm}.
\item Laggard coefficient: \blank{1cm} (negative/positive).
\item Leader coefficient: \blank{1cm} (sign).
\end{enumerate}

\section*{Round 3: Level 2 Prompts}
\begin{enumerate}
\item Explain the non-monotonic pattern across environmental bins.
\item Why is ``smart money'' consistent with risk management without moral endorsement?
\item Write the empirical spec.
\item Contrast FSU with Hong-Kacperczyk (2009) on sin-stock exclusion.
\end{enumerate}

\section*{Round 4: Minimal Scaffolding}
From a blank page: (i) data, (ii) binning, (iii) results, (iv) interpretation, (v) two caveats.

\section*{Round 5: Blank Page}
Essay: what does institutional ESG behaviour reveal about corporate environmental policy?

\vspace{6pt}
\begin{drillbox}
\textbf{Speed Drill.} (1) Data years. (2) Three environmental bins. (3) Laggard holdings sign. (4) Leader holdings sign. (5) Smart-money reading.
\end{drillbox}

\section*{Quick Reference \& Answer Key}
\begin{answerbox}
\begin{itemize}
\item \textbf{Data.} KLD 1992--2008, 13-F holdings.
\item \textbf{Bins.} Laggard / neutral / leader.
\item \textbf{Result.} Both laggards and leaders are underweighted; neutral is preferred.
\item \textbf{Interpretation.} Risk-management, not moral endorsement.
\end{itemize}
\end{answerbox}

\begin{keybox}
\textbf{30-second exam pitch.} Fernando, Sharfman \& Uysal (2017) challenge the simple ``institutions love green'' narrative. Using KLD environmental data 1992--2008 and Thomson 13-F institutional holdings, they classify firms as laggards, neutral, or leaders based on environmental strengths minus concerns. Institutional ownership is highest for \emph{neutral} firms: laggards are underweighted because of environmental tail-risk, and leaders are also underweighted (more mildly), likely reflecting concerns that aggressive environmental investment may be inefficient or destroy shareholder value. The pattern is strongest for long-horizon institutions (pensions). The paper reframes institutional ESG behaviour as risk management rather than moral endorsement --- institutions avoid environmental costs without necessarily paying a premium for best-in-class green performance.
\end{keybox}

\end{document}
